% Section 1: Opening Movement
%
% Regulator-facing statement of the problem, lens, and evidentiary taxonomy.

\section{Opening Movement}
\label{sec:opening}

The Financial Services AI Risk Management Framework (FS~AI~RMF), developed by the Cyber Risk Institute with the Financial Services Sector Coordinating Council and the US Treasury and released in February 2026, specifies 230 control objectives spanning governance, data, model development, validation, monitoring, third-party risk, and consumer protection \parencite{treasuryFSAIRMFRelease2026}. Read as a checklist, the framework is operationally daunting and conceptually opaque: community banks still face the question of what \emph{commensurate} compliance means, and the document itself does not answer it. The FS~AI~RMF, unlike the supervisory model-risk guidance issued the same year (SR Letter~26-2), has no asset-size threshold; it is designed to apply to all financial institutions regardless of size, type, complexity, or criticality. Read as a structure, however, the framework admits a generative organizing lens that makes interpretation tractable and gives operationalization a structure to build on: a wide range of control objectives can be understood as protecting an interest of a party who is not in the room when an AI system makes or supports a financial decision. We develop empty-chair representation as an interpretive lens, demonstrated through worked examples across control families rather than asserted as a coding of all 230 objectives; whether the reading holds framework-wide is an empirical question we frame, but do not settle, here (a systematic content-analytic coding with trained raters is future work).

Recent commentary on the FS~AI~RMF describes it as a \emph{systems blueprint} rather than a checklist, and urges institutions to map controls to architecture, ownership, and evidence rather than to documentation alone~\parencite{mushahwar2026}. We extend that move one step further: the blueprint has a normative structure, and that structure is \emph{empty-chair representation}: the design of AI governance such that specific absent parties are structurally represented in the architecture of decision-making, not merely in its documentation.

The party in the empty chair varies by control. In some, it is the loan applicant who will be denied without recourse. In others, the depositor whose funds are deployed in ways they did not consent to. In others, the future compliance officer who will inherit responsibility for decisions made under regimes they did not construct. In others, the examiner applying updated model-risk guidance that had not yet been issued when the system was deployed. In others, the institution itself in its future state, after personnel turnover and ownership changes. The empty chair is plural, and the chairs do not always agree. The proposed architectural function of regulatory governance is to make these absent interests inspectable by the parties present at the moment of decision and by those who later review it.

This reframe matters operationally. By \emph{verification regime} we mean the set of technical, procedural, and evidentiary mechanisms by which a party not present at the moment of decision can later test a material claim about that decision, rather than accept the claim from the outcome or from an institution's assertion. A control \emph{represents} an absent party when it induces observable, auditable constraints on system behavior or evidence production that a later examiner can use to assess whether that party's interests were respected. Representation in this sense is architectural; documentation can participate in it, but an assertion of representation is not evidence that representation occurred.

The object being verified must be named. A decision record may preserve the \emph{inputs and context} available at decision time. SHAP can characterize \emph{functional feature dependence} under a specified value function and background construction; LIME can provide a \emph{local behavioral approximation} under a specified neighborhood and surrogate \parencite{lundberg2017,ribeiro2016}. Attention and activation analyses expose selected \emph{internal computational states}, although attention weight alone does not establish feature importance \parencite{jain2019}. Chain-of-thought and other natural-language accounts are \emph{generated rationales}. A contemporaneous reviewer record is evidence about a \emph{human deliberative record}: what was presented, consulted, recorded, and decided, not transparent access to every cognitive process. A signature or approval log records \emph{institutional authorization and review}. These artifacts can each be useful. They are not interchangeable, and an artifact may faithfully establish one object without establishing another.

This distinction changes the examination question. The question is not simply whether the institution has produced an explanation. It is \emph{what evidentiary object does this artifact expose, what conclusion is the recipient being asked to draw from it, and what binds the artifact to that conclusion?} A local approximation may illuminate behavior without recording deliberation; an understandable rationale may communicate a decision without faithfully characterizing functional dependence; a signature may establish authorization without establishing the truth of the material reviewed. Verification fails when the claimed conclusion requires a different object from the one the architecture exposes.

Examiners reviewing AI deployments at community banks face a problem the framework does not solve for them: 230 control objectives do not themselves specify how documentary evidence warrants a conclusion about deployed behavior. With an empty-chair frame, examination becomes a structured inquiry: \emph{which absent parties does this AI system affect, what artifact represents their interests, and what conclusion can that artifact support?} The frame supplements the checklist with an evidentiary question.

Two positions follow. First, \emph{an explanation artifact is not, by itself, a verification regime for a different evidentiary object}. Post-hoc methods have legitimate uses, including model diagnosis, communication, and exploration; their limits depend on their target and construction \parencite{oh2024}. But in adversarial review the recipient cannot assume that an understandable artifact establishes a different proposition without an inspectable binding between them \parencite{bordt2022}. SHAP and LIME can be manipulated in ways that conceal objectionable model behavior \parencite{slack2020}; generated chain-of-thought can be unfaithful to the computation that produced an answer \parencite{lanham2023,turpin2023}. These findings concern different artifacts and mechanisms. Their common governance lesson is narrower: the institution must establish the warrant for the inference it asks the examiner or affected party to draw.

Second, \emph{review of an explanation verifies that the explanation was reviewed; it does not by itself verify a decision process that the explanation does not faithfully represent}. A signature can establish authorization, responsibility, and completion of a specified procedure. Those are real governance functions. If the institution also claims that the signature establishes the substantive basis of an AI recommendation, it must show what the reviewer could inspect, what independent judgment the procedure required, and how the reviewed artifact bears on that substantive claim. Human presence in a workflow and verification of a model-supported decision are not equivalent propositions.

These positions do not counsel against AI in banking. They distinguish uses by the evidentiary claim the architecture must support. Pattern recognition, anomaly detection, and consistency monitoring can surface material for investigation while leaving the disposition and its documented basis with an authorized human. Decision-recommendation systems may demand evidence about model behavior, data, policy fit, human review, and the relation among them. The difference is not captured by a generic ``human in the loop'' label or a generic explanation requirement; it depends on which party decides, which artifact is exposed, and which conclusion that artifact can support.

A note on what the empty-chair frame is doing structurally, before the body sections develop it. Absence in regulated decision contexts is not always a static gap. An interface may make challenge costly; a documentation regime may preserve only later justification; a metric may omit the evidence needed to test the conclusion drawn from it. The frame therefore asks not only \emph{who is absent} but \emph{what access does the architecture remove, what artifact replaces it, and what inference is drawn from the replacement?} Section~\ref{sec:diagnostics} develops this recurring inference gap as silence-manufacture. The claim is conditional: silence is not evidence of absence when the system being evaluated materially contributes to that silence.

The remainder of this paper demonstrates the lens against specific FS~AI~RMF control objectives, develops three diagnostic questions, and derives four provisional architectural requirements at the property level. Its test is practical: whether the lens helps a community bank or examiner distinguish the evidentiary object an artifact exposes from the broader conclusion the institution asks that artifact to support.
